% Intervention algebra. One claim: the five causal edits compose into a single
% compiled intervention, and the dangerous one (erase) only counts once it carries
% a receipt: recovery driven to chance is CALIBRATED, a sham that leaves it at 1.0
% stays EXPLORATORY. Causal tier, so amber (rlTierCausal) throughout.
\documentclass[border=8pt]{standalone}
% House preamble for every reward-lens TikZ figure.
% The figure file sets \documentclass{standalone} and, before it is \input, the build
% sets \rltheme (light|dark). This file loads the matching palette and defines the shared
% visual vocabulary, so consistency across figures comes from here, not from discipline.
%
% Engine: lualatex or xelatex gives the real Inter (matches the site). pdflatex falls back
% to Fira Sans, which is close. Either way the figures do not speak Computer Modern.

\usepackage{iftex}
\ifLuaTeX
  \usepackage{fontspec}\setsansfont{Inter}[BoldFont={Inter SemiBold}, ItalicFont={Inter Italic}]
\else\ifXeTeX
  \usepackage{fontspec}\setsansfont{Inter}[BoldFont={Inter SemiBold}, ItalicFont={Inter Italic}]
\else
  \usepackage[T1]{fontenc}\usepackage[sfdefault]{FiraSans}
\fi\fi
\renewcommand{\familydefault}{\sfdefault}

\usepackage{amsmath}
\usepackage{tikz}
\usetikzlibrary{
  positioning, calc, arrows.meta, fit, backgrounds, shapes.geometric,
  shapes.misc, decorations.pathreplacing, decorations.pathmorphing,
  shadows.blur, patterns, math
}

% AUTO-GENERATED from palette.json by emit_palette.py. Do not edit by hand.
% Each figure sets \rltheme to light or dark before \input; the preamble loads this file.
\usepackage{etoolbox}
\providecommand{\rltheme}{light}%
\ifdefstring{\rltheme}{dark}{%
  \definecolor{rlInk}{HTML}{E6E8EB}%
  \definecolor{rlInkTwo}{HTML}{B7BDC7}%
  \definecolor{rlMuted}{HTML}{9AA0AA}%
  \definecolor{rlLine}{HTML}{3A3D44}%
  \definecolor{rlSurface}{HTML}{1D1E22}%
  \definecolor{rlSurfaceTwo}{HTML}{26272C}%
  \definecolor{rlBaseline}{HTML}{55565E}%
  \definecolor{rlTierObs}{HTML}{2DD4BF}%
  \definecolor{rlTierCausal}{HTML}{F59E0B}%
  \definecolor{rlTierVuln}{HTML}{FB7185}%
  \definecolor{rlTrustExp}{HTML}{64748B}%
  \definecolor{rlTrustCal}{HTML}{60A5FA}%
  \definecolor{rlTrustReg}{HTML}{A78BFA}%
  \definecolor{rlTrustAdj}{HTML}{22C55E}%
  \definecolor{rlSignal}{HTML}{E879F9}%
  \definecolor{rlGauge}{HTML}{22D3EE}%
  \definecolor{rlChosen}{HTML}{5598E7}%
  \definecolor{rlRejected}{HTML}{F0716F}%
  \definecolor{rlSeriesB}{HTML}{F98A5C}%
  \definecolor{rlAccent}{HTML}{FB8C5A}%
  \definecolor{rlSeqLo}{HTML}{0D366B}%
  \definecolor{rlSeqHi}{HTML}{CDE2FB}%
  \definecolor{rlDivMid}{HTML}{2A2B31}%
}{%
  \definecolor{rlInk}{HTML}{16181D}%
  \definecolor{rlInkTwo}{HTML}{4B5563}%
  \definecolor{rlMuted}{HTML}{6B7280}%
  \definecolor{rlLine}{HTML}{D7DAE0}%
  \definecolor{rlSurface}{HTML}{FFFFFF}%
  \definecolor{rlSurfaceTwo}{HTML}{F4F5F7}%
  \definecolor{rlBaseline}{HTML}{C3C2B7}%
  \definecolor{rlTierObs}{HTML}{0D9488}%
  \definecolor{rlTierCausal}{HTML}{B45309}%
  \definecolor{rlTierVuln}{HTML}{BE123C}%
  \definecolor{rlTrustExp}{HTML}{94A3B8}%
  \definecolor{rlTrustCal}{HTML}{3B82F6}%
  \definecolor{rlTrustReg}{HTML}{7C3AED}%
  \definecolor{rlTrustAdj}{HTML}{15803D}%
  \definecolor{rlSignal}{HTML}{A21CAF}%
  \definecolor{rlGauge}{HTML}{0891B2}%
  \definecolor{rlChosen}{HTML}{2A78D6}%
  \definecolor{rlRejected}{HTML}{E34948}%
  \definecolor{rlSeriesB}{HTML}{EB6834}%
  \definecolor{rlAccent}{HTML}{E85D2C}%
  \definecolor{rlSeqLo}{HTML}{CDE2FB}%
  \definecolor{rlSeqHi}{HTML}{0D366B}%
  \definecolor{rlDivMid}{HTML}{F0EFEC}%
}%
 % defines the rl* colors for the current \rltheme

% digit-name aliases, so both rlInkTwo and rlInk2 resolve
\colorlet{rlInk2}{rlInkTwo}
\colorlet{rlSurface2}{rlSurfaceTwo}

% ---- global defaults -------------------------------------------------------
\tikzset{
  >={Stealth[length=2.4mm, width=1.8mm]},
  font=\sffamily,
  every node/.append style={text=rlInk},
}

% ---- chips: the three-chip vocabulary (tier / gauge status / works-on) ------
% Outline chip reads correctly on both themes because the role color is a mid-tone.
\tikzset{
  rlchip/.style 2 args={
    draw=#1, text=#1, line width=0.5pt, rounded corners=2.5pt,
    inner xsep=5pt, inner ysep=2.6pt, font=\sffamily\bfseries\scriptsize,
    label/.expanded=#2,
  },
  rlchip/.default={rlMuted}{},
  % filled variant, surface-colored text so it contrasts in both themes
  rlchipfill/.style={
    draw=#1, fill=#1, text=rlSurface, rounded corners=2.5pt,
    inner xsep=5pt, inner ysep=2.6pt, font=\sffamily\bfseries\scriptsize,
  },
}
% simple outline chip taking one color (most common)
\newcommand{\rlchipnode}[3][rlMuted]{\node[draw=#1, text=#1, line width=0.6pt,
  rounded corners=2.5pt, inner xsep=5pt, inner ysep=2.6pt,
  font=\sffamily\bfseries\scriptsize] (#2) {#3};}

% ---- cards and panels ------------------------------------------------------
\tikzset{
  rlcard/.style={
    draw=rlLine, fill=rlSurface, line width=0.7pt, rounded corners=4pt,
    blur shadow={shadow blur steps=7, shadow xshift=0pt, shadow yshift=-1pt,
      shadow opacity=18, shadow blur radius=3pt},
  },
  rlpanel/.style={draw=rlLine, fill=rlSurfaceTwo, line width=0.6pt, rounded corners=3pt},
  rlfield/.style={font=\sffamily\scriptsize, text=rlMuted, align=left},
  rlvalue/.style={font=\sffamily\bfseries, text=rlInk},
  rllabel/.style={font=\sffamily\footnotesize, text=rlInk2},
  rlnote/.style={font=\sffamily\scriptsize\itshape, text=rlMuted},
}

% ---- lines, axes, vectors --------------------------------------------------
\tikzset{
  rlaxis/.style={draw=rlLine, line width=0.7pt, -{Stealth[length=1.8mm]}},
  rlbaseline/.style={draw=rlBaseline, line width=0.6pt, dash pattern=on 2pt off 2pt},
  rlflow/.style={draw=rlInk2, line width=0.9pt, -{Stealth[length=2.4mm]}},
  rlvec/.style={line width=1.3pt, -{Stealth[length=2.8mm, width=2.2mm]}},
  rldrop/.style={draw=rlMuted, line width=0.5pt, dash pattern=on 1.6pt off 1.6pt},
}

% ---- gate glyph: a narrow keyhole the evidence must pass to climb a rung ----
% \rlgate{name}{fill color}{label}  draws a small rounded gate node.
\newcommand{\rlgate}[3]{%
  \node[draw=#2, fill=rlSurface, line width=0.9pt, rounded corners=2pt,
        minimum width=6.5mm, minimum height=6.5mm, inner sep=0pt] (#1) {};
  \node[text=#2, font=\sffamily\bfseries\footnotesize] at (#1) {#3};%
}

% ---- answer-key glyph (calibration): a tiny key ----------------------------
\newcommand{\rlkey}[2][rlTrustAdj]{%
  \begin{scope}[shift={(#2)}, rotate=45]
    \draw[fill=#1, draw=#1] (0,0) circle (1.6mm);
    \fill[rlSurface] (0,0) circle (0.6mm);
    \draw[#1, line width=1.1pt] (0,-1.6mm) -- (0,-4.4mm);
    \draw[#1, line width=1.1pt] (0,-3.0mm) -- (1.1mm,-3.0mm);
    \draw[#1, line width=1.1pt] (0,-4.0mm) -- (0.9mm,-4.0mm);
  \end{scope}%
}

% ---- frozen-study seal: a small notched stamp ------------------------------
\newcommand{\rlseal}[2][rlTrustReg]{%
  \node[star, star points=12, star point ratio=0.9, draw=#1, fill=#1,
        text=rlSurface, minimum size=9mm, inner sep=0pt,
        font=\sffamily\bfseries\tiny] at (#2) {\#hash};%
}

\begin{document}
\begin{tikzpicture}[scale=1.0]

  % ---- panel -------------------------------------------------------------
  \fill[rlSurfaceTwo, rounded corners=6pt] (-0.65,-0.7) rectangle (14.5,3.55);

  % ---- one compiled intervention: brace + label over the strip -----------
  \draw[decorate, decoration={brace, amplitude=5pt}, line width=0.8pt, rlMuted]
    (0.45,2.44) -- (11.95,2.44);
  \node[font=\sffamily\bfseries\small, text=rlInk] at (6.2,3.2)
    {one compiled intervention};
  \node[font=\sffamily\scriptsize, text=rlMuted] at (6.2,2.86)
    {{\ttfamily compose(\ldots)} mounts them in declaration order at one shared site $\ell$};

  % ---- the five interlocking blocks --------------------------------------
  % Each block is a jigsaw tile: a socket on the left, a knob on the right, so
  % consecutive blocks snap along a shared wavy seam.
  \newcommand{\ivblock}[5]{% #1 cx  #2 name  #3 gloss  #4 leftsocket(0/1)  #5 rightknob(0/1)
    \pgfmathsetmacro{\cx}{#1}%
    \def\cy{1.6}\def\hw{1.15}\def\hh{0.72}\def\kr{0.26}%
    \draw[line width=0.9pt, draw=rlTierCausal, fill=rlSurface]
      (\cx-\hw,\cy-\hh) -- (\cx+\hw,\cy-\hh)
      \ifnum#5=1 -- (\cx+\hw,\cy-\kr) arc[start angle=-90, end angle=90, radius=0.26cm] -- (\cx+\hw,\cy+\hh)
      \else -- (\cx+\hw,\cy+\hh)\fi
      -- (\cx-\hw,\cy+\hh)
      \ifnum#4=1 -- (\cx-\hw,\cy+\kr) arc[start angle=90, end angle=-90, radius=0.26cm] -- (\cx-\hw,\cy-\hh)
      \else -- (\cx-\hw,\cy-\hh)\fi
      -- cycle;
    \node[font=\sffamily\bfseries\small, text=rlTierCausal] at (\cx,\cy+0.18) {#2};
    \node[font=\sffamily\scriptsize, text=rlMuted] at (\cx,\cy-0.24) {#3};
  }
  \ivblock{1.6}{patch}{swap in a value}{0}{1}
  \ivblock{3.9}{steer}{add a direction}{1}{1}
  \ivblock{6.2}{ablate}{zero a subspace}{1}{1}
  \ivblock{8.5}{edit}{rewrite a value}{1}{1}
  \ivblock{10.8}{erase}{remove all trace}{1}{0}

  % ---- receipt stamped onto the erase block ------------------------------
  % pin from the erase block out to the certificate
  \fill[rlTierCausal] (11.95,1.5) circle (1.1pt);
  \draw[rlTierCausal, line width=0.9pt] (11.95,1.5) -- (12.42,1.42);

  % the certificate card (a stamped, perforated receipt), in the calibrated blue
  \def\rx{12.42}\def\rX{14.2}\def\ry{0.5}\def\rY{2.18}
  \draw[rlTrustCal, line width=1.0pt, fill=rlSurface,
        blur shadow={shadow blur steps=7, shadow xshift=0.4pt, shadow yshift=-1.2pt,
                     shadow opacity=22, shadow blur radius=3pt}]
    (\rx,\ry) rectangle (\rX,\rY);
  % perforation strip near the top (tear-off holes)
  \foreach \px in {12.6,12.8,...,14.0}{\fill[rlSurfaceTwo] (\px,1.98) circle (0.5mm);}
  \node[font=\sffamily\scriptsize\itshape, text=rlMuted] at (13.31,2.06) {erasure certificate};
  \node[font=\sffamily\scriptsize, text=rlInk2] at (13.31,1.6) {recovery AUC};
  \node[font=\sffamily\bfseries\large, text=rlTrustCal] at (13.31,1.26) {0.51};
  \node[draw=rlTrustCal, fill=rlTrustCal, text=rlSurface, rounded corners=2.5pt,
        inner xsep=5pt, inner ysep=2.6pt, font=\sffamily\bfseries\scriptsize]
    at (13.31,0.84) {CALIBRATED};

  % the sham counterfactual: no drop, no receipt, no lift off the floor
  \node[font=\sffamily\scriptsize, text=rlMuted, align=center] at (13.31,0.06)
    {sham keeps AUC 1.00\\[-1pt]{\color{rlTrustExp}stays exploratory}};

  % ---- the claim, in plain words -----------------------------------------
  \node[anchor=north west, align=left, font=\sffamily\footnotesize\itshape, text=rlMuted]
    at (-0.5,0.28)
    {\texttt{compose()} chains five causal edits into one hook at one site.\\
     Erase is the dangerous one: it must show a receipt, or it stays exploratory.};
\end{tikzpicture}
\end{document}
